\documentclass[11pt]{article}
\usepackage[margin=1in]{geometry}
\usepackage[T1]{fontenc}
\usepackage{lmodern}
\usepackage{amsmath,amssymb,amsthm,booktabs,array,microtype}
\usepackage[colorlinks=true,linkcolor=blue,citecolor=blue,urlcolor=blue]{hyperref}
\usepackage{enumitem}
\setlist{nosep}
\setlength{\parskip}{4pt}
\newtheorem{theorem}{Theorem}
\newtheorem{lemma}[theorem]{Lemma}
\newtheorem{corollary}[theorem]{Corollary}
\newtheorem{proposition}[theorem]{Proposition}
\newcommand{\Z}{\mathbb Z}
\newcommand{\ind}{\mathbf 1}
\newcommand{\U}{\mathcal U}
\DeclareMathOperator{\lcm}{lcm}
\title{A Uniform Optimum for Short-Interval\protect\\Square-Difference Masks\protect\\[5pt]\large A computer-assisted result for a 256-bit raw-table model}
\author{GPT-6 astra}
\date{10 October 2026}
\begin{document}
\maketitle
\begin{abstract}
Fermat's factorization method searches for a representation of an integer
$N$ as a difference of two squares, $N=a^2-b^2$. A modular filter can reject
a candidate $a$ whenever $a^2-N$ is not a square modulo a chosen integer
$m$. The filter repeats every $m$ positions, so its decisions can be stored
in a table with one bit per residue. Which modulus gives the strongest
guaranteed filtering when this table is limited to 256 bits?

We measure filtering strength by the largest number of candidates that can
pass in any $h$ consecutive positions, taking the worst case over interval
starts and inputs coprime to the chosen modulus. We prove that $m=144$
minimizes this worst-case count among every integer modulus
$1\le m\le256$, simultaneously for every positive interval length $h$.
For example, at most two candidates pass in any interval of 18 positions,
and at most four pass in any interval of 48 positions. An explicit formula
for the count at modulus 144, independently checked finite computations,
and a density comparison together prove the statement for all lengths.

The theorem concerns candidate counts under this specific table-storage
model; it does not establish a faster factorization algorithm. Earlier
work already studies modular filtering and uses modulus 144. Whether the
precise all-length optimality theorem is new remains unresolved.
\end{abstract}

\section{Definition and scope}

Write $Q_m=\{b^2\bmod m:b\in\Z\}$ and let $\U_m$ denote the unit residues
modulo $m$. Set $\U_1=\{0\}$ and $Q_1=\{0\}$. Define
\[
 S_m(n)=\{a\bmod m:a^2-n\bmod m\in Q_m\},\qquad
 C_{m,n,s}(h)=\sum_{j=0}^{h-1}\ind_{S_m(n)}(s+j),
\]
\[
 F_m(h)=\max_{n\in\U_m,\;0\le s<m}C_{m,n,s}(h),\qquad
 G_B(h)=\min_{1\le m\le B}F_m(h).
\]
Positions are counted with multiplicity when $h>m$. We also use
$C_{m,n,s}(0)=F_m(0)=0$.

The storage charge is exactly $m$ logical bits for a one-bit-per-residue
table for the current $n$. The limit $B=256$ therefore represents 32 bytes
of raw table payload. Metadata, construction workspace, verifier memory,
and a cache over all values of $n$ are excluded. This is an explicit
uncompressed representation. In fact, every even-$m$ mask has period
dividing $m/2$, and parity permits further packing in some cases. No claim
of optimal compressed storage follows.

The maxima over $n$ and $s$ give uniform bounds for arbitrary intervals.
They need not be attained at the correlated start and length of an actual
Fermat factorization search. The verification uses finite residue sets,
without accepting an integer to factor or testing any deployed key.

\begin{theorem}\label{thm:main}
For every integer $h\ge1$ and every integer $1\le m\le256$,
\[
                         F_{144}(h)\le F_m(h).
\]
Thus $G_{256}(h)=F_{144}(h)$ for all positive lengths. The assertion allows
ties at individual lengths and does not extend the modulus budget.
\end{theorem}

The proof separates an analytic description of one periodic family from
finite, computer-assisted comparisons. The finite lemmas are not claimed
to have been formalized in a proof assistant. The package includes two
independent implementations, attaining witnesses, exact output hashes,
and the verification limits needed to audit those lemmas.

\section{An exact formula for the modulus 144}

Define two functions on nonnegative integers by
\[
 B(h)=2\left\lfloor\frac h{24}\right\rfloor+b(h\bmod24),\qquad
 Q(h)=4\left\lfloor\frac h{72}\right\rfloor+k(h\bmod72),
\]
where
\[
b(r)=\begin{cases}0&r=0,\\1&1\le r\le6,\\2&7\le r\le23,\end{cases}
\qquad
k(r)=\begin{cases}
0&r=0,\\1&1\le r\le2,\\2&3\le r\le18,\\
3&19\le r\le26,\\4&27\le r\le71.
\end{cases}
\]

\begin{lemma}\label{lem:envelope}
For every integer $h\ge0$, $F_{144}(h)=\max\{B(h),Q(h)\}$.
\end{lemma}
\begin{proof}
CRT separates the conditions modulo $16$ and $9$. The square residues
modulo $16$ are $0,1,4,9$. Direct substitution gives the complete unit table
\[
\begin{array}{c|c}
n\pmod {16}&\text{allowed }a\\ \hline
1,13&1,7\pmod8\\
5,9&3,5\pmod8\\
3,11&2\pmod4\\
7,15&0\pmod4
\end{array}
\]
Modulo $9$, $n=2,5,8$ gives precisely $a=0\pmod3$. The values $n=1,4,7$
give respectively $a=\pm u\pmod9$ with $u=1,2,4$.

First combine $a=0\pmod3$ with the modulo-$16$ patterns. A single class
modulo $4$ becomes one class modulo $12$. A two-class pattern modulo $8$
becomes two points, separated by $6$, in a period of $24$. The latter
family has envelope $B(h)$ and dominates the former.

Next combine the two-root pattern modulo $9$ with a single class modulo
$4$. There are two points per period $36$, with shorter cyclic separation
$4$, $8$, or $16$. The largest envelope is obtained with separation $4$.
Its thresholds for one through four points within a period of $72$ are
$1,5,37,41$.

It remains to combine two roots modulo $9$ with two classes modulo $8$.
Translation by $36$ exchanges the two possible modulo-$8$ patterns and
preserves the modulo-$9$ roots. It suffices to intersect $a=\pm1\pmod8$
with $a=\pm u\pmod9$. The resulting masks are
\[
\begin{array}{c|l|l|l}
u&\text{positions in }[0,71]&\text{cyclic gaps}&
 \text{least lengths for 1,2,3,4 points}\\ \hline
1&1,17,55,71&16,38,16,2&1,3,19,35\\
2&7,25,47,65&18,22,18,14&1,15,33,51\\
4&23,31,41,49&8,10,8,46&1,9,19,27
\end{array}
\]
A window containing successive selected points requires one plus the sum
of the intervening gaps. Maximizing these three profiles gives $Q(h)$.
The preceding single-class-modulo-$4$ case is dominated, since its
thresholds $1,5,37,41$ are no earlier than $1,3,19,27$ and both have four
points per period $72$. All CRT combinations used here are realized by
unit residues $n$, so the maximum of the remaining two envelopes is exact.
\end{proof}

\begin{corollary}\label{cor:error}
For every integer $h\ge0$,
\[
                         F_{144}(h)\le\frac h{12}+\frac74.
\]
The additive constant $7/4$ is attained.
\end{corollary}
\begin{proof}
In $B(h)-h/12$, complete $24$-blocks cancel. The largest remainder excess
is $2-7/12=17/12$. In $Q(h)-h/12$, each complete $72$-block reduces the
excess by $2$. The remainder jumps at $1,3,19,27$, where the excesses are
$11/12,7/4,17/12,7/4$. Thus $7/4$ is an upper bound. For $n=1$ and the
three-position interval $-1,0,1$, exactly two positions pass, attaining
$2-3/12=7/4$. A second example is $n=97$, interval $23,\ldots,49$, with
four passing positions and excess $4-27/12=7/4$.
\end{proof}

\section{Finite comparisons and the infinite tail}

Put $R_m(n)=|S_m(n)|$ and $\rho_m=\max_{n\in\U_m}R_m(n)/m$.
\begin{lemma}\label{lem:average}
For all $m\ge1$ and $h\ge0$, $F_m(h)\ge\lceil h\rho_m\rceil$.
\end{lemma}
\begin{proof}
Choose a mask of maximum cardinality. In the sum of its window counts over
the $m$ starts, every marked residue occurs exactly $h$ times. Its average
window count is therefore $h\rho_m$; the maximum is at least its ceiling.
\end{proof}

\begin{lemma}[Finite certificate]\label{lem:finite}
Complete enumeration in the stated finite domain establishes:
\begin{enumerate}[label=(\roman*)]
\item $\rho_{144}=1/12$ and $\rho_m\ge2/21$ for every
      $1\le m\le256$, $m\ne144$. The latter minimum occurs at $252$.
\item $F_{144}(h)\le F_m(h)$ for every $1\le h\le146$ and
      $1\le m\le256$.
\end{enumerate}
\end{lemma}

The computational evidence for this lemma is described explicitly in
Section~\ref{sec:certificate}. As a compact check of (i), the complete list
of moduli in this domain with $\rho_m<1/8$ is
\[
\begin{array}{c|cccc}m&144&180&240&252\\\hline
\rho_m&1/12&1/10&1/10&2/21\end{array}.
\]
The proof does not assume that low density by itself orders short windows.

\begin{proof}[Proof of Theorem~\ref{thm:main}]
Lemma~\ref{lem:finite}(ii) proves the claim through $h=146$. For a competing
modulus and $h\ge147$, Lemmas~\ref{lem:average} and~\ref{lem:finite},
together with Corollary~\ref{cor:error}, give
\[
 F_m(h)\ \ge\ \frac{2h}{21}
 \ \ge\ \frac h{12}+\frac74
 \ \ge\ F_{144}(h).
\]
The middle inequality is exactly $h/84\ge7/4$, or $h\ge147$.
This covers every positive integer length. For $h>147$ the middle
inequality is strict, proving unique optimality at each $h\ge148$ as an
additional sufficient conclusion.
\end{proof}

The original, prospectively fixed horizon was $1\le h\le64$. The all-length
extension was recorded after the common winner and density gap were
observed. The original results were preserved. The extension changed
neither the mask definition nor the budget or admissible moduli.

\section{What the optimum says about short windows}

The exact envelope yields the following original-horizon values:
\[
\begin{array}{c|rrrrrr}
h&1\text{--}2&3\text{--}18&19\text{--}26&27\text{--}48&49\text{--}54&55\text{--}64\\\hline
G_{256}(h)&1&2&3&4&5&6
\end{array}.
\]
The least raw-bit budgets attaining these final-budget optima are
\[
\begin{array}{c|rrrrrrrr}
h&1&2&3&4&5\text{--}6&7\text{--}8&9\text{--}12&13\text{--}64\\\hline
\text{least budget}&1&4&3&4&9&16&36&144
\end{array}.
\]
Individual-length ties are retained in the certificate. At $h=20$ the
optimal moduli are $144,180$; at $h=32$ they are $144,180,252$. Thus the
common optimum need not be the cheapest choice for a particular short
length.

There is also an elementary budget-independent consequence. For any
modulus and $h\ge3$, choosing $n=1$ admits both $a=-1$ and $a=1$ in a
three-position interval, so $F_m(h)\ge2$. The exact $144$ envelope attains
two for $3\le h\le18$. Consequently for \emph{every} budget $B\ge144$,
$G_B(h)=2$ throughout that range, even though the all-length optimality
theorem fixes $B=256$.

Density ordering is insufficient as a general short-window rule. The
worst densities satisfy $\rho_9=1/3<1/2=\rho_4$, but
$F_9(2)=2>1=F_4(2)$. Modulo $9$, $n=7$ admits adjacent positions $4,5$.
Also, keeping only maximum-density masks is invalid: the maximum-density
$n=2,5,8$ masks are $\{0,3,6\}$ and have two-window maximum one, whereas
the lower-density $n=7$ mask has maximum two. Nevertheless, in the stated
global budget problem the density minimizer $144$ is also the uniform
window optimum. The experiment did not refute that global selection rule
in this domain.

\section{The finite certificate and its verification}\label{sec:certificate}

The first construction enumerates every $m=1,\ldots,256$, every unit $n$,
and every start. After forming the exact mask, it updates all window sums
by the recurrence
\[
 W_s(0)=0,\qquad W_s(h+1)=W_s(h)+f_{m,n}(s+h).
\]
Taking the maximum over all enumerated $n,s$ gives both an upper bound and
an attaining witness for each $F_m(h)$. There is no sampling, numerical
approximation, or squareclass pruning. A witness alone certifies only a
lower bound; the full enumeration supplies the upper bound.

The independent implementation constructs masks from the equation
$n=a^2-b^2\pmod m$, then uses repeated integer bitsets and population
counts for each window. It agrees on all 16,384 original table entries,
all original attaining witnesses, all budget/length tie sets, and all
maximum-density residue ties.

\begin{center}
\begin{tabular}{lrr}
\toprule
Construction coverage&$h\le64$&$h\le146$\\
\midrule
Unit masks&19,948&19,948\\
Mask predicates&3,406,801&3,406,801\\
Window updates&218,035,264&497,392,946\\
Attaining witnesses&16,384&37,376\\
Direct witness bit checks&532,480&2,747,136\\
\bottomrule
\end{tabular}
\end{center}

The second implementation does not independently recompute every competing
upper bound beyond $64$. Instead it directly checks the analytic $144$
envelope through $146$ and certifies all 20,910 competing cells for
$65\le h\le146$ by
$\lceil h\rho_m\rceil\ge F_{144}(h)$. This suffices for optimality.
It also directly checks all 37,376 extended attaining witnesses. An
additional constants checker enumerates 995,328 remainder cases for
$m=144$ and confirms the excess $7/4$ and the density gap.

All arithmetic used for correctness is integral or rational. Each exhaustive
verifier and certificate checker accepts no arguments and enforces limits
of 120 CPU seconds and 256 MiB
of address space. Those limits govern the verifier, not the logical mask
payload. Mathematical JSON is deterministic; runtime telemetry is stored
separately. Parent reruns and document-build checks are recorded in
\texttt{parent\_verification.json} and \texttt{SHA256SUMS}.

To audit the central result on Linux with Python 3.10 or later, run the
following from the extracted package directory. No third-party Python
packages are required.
\begin{quote}\small\ttfamily
sha256sum -c SHA256SUMS\\
python3 wave2\_b/test\_minimax\_146.py\\
python3 wave2\_b/generate\_minimax\_146.py\\
python3 wave2\_b/check\_all\_h\_constants.py\\
python3 wave2\_c/verify\_minimax.py\\
python3 wave2\_c/verify\_peer.py\\
python3 wave2\_c/extend\_optimality.py\\
python3 wave2\_c/verify\_extended\_peer.py
\end{quote}
The original $64$-length construction can also be regenerated with
\texttt{wave2\_b/generate\_minimax.py}. Reruns replace their result and
telemetry files; preserve the supplied record by using a separate copy.
The mathematical hashes are expected to agree. Hash agreement identifies
artifacts and does not substitute for proof review.

\section{Prior art, novelty, and release decision}

Square-difference filters and modular-hyperbola projections have substantial
prior art. Bower, Evans, Luo, and Miller give CRT and prime-power
projection counts~\cite{bower}. Hittmeir develops the hyperbolic sieve,
prime-power counts, modulus-choice examples, and a stronger growing-modulus
factorization analysis~\cite{hittmeir}. Those results already cover the
underlying full-period mathematics.

Hewitt's 2022 empirical study explicitly searches 792 Fermat sieve moduli
formed from bounded powers of $2,3,5,7$. It constructs a Pareto front for
full-period reduction versus compact passing-entry table size, conditioned
on 24 input categories~\cite{hewitt}.
Its tables include 144. Thus neither finite modulus optimization nor the
choice of 144 is new here. The exact worst-over-unit-and-start,
all-length statement under the raw 256-bit model was not located in the
bounded search. That absence does not establish historical priority.

This research note and its accompanying evidence package are intended for
public reproducibility and further mathematical review. The result is a
verified finite-budget combinatorial theorem whose historical novelty
remains unestablished. It establishes candidate-count optimality under the
stated representation model, without establishing a faster factorization
algorithm. Its release classification is \texttt{PUBLIC\_RESEARCH}.

A further mathematical study should explicitly change one assumption at
a time, for example the representation cost or modulus budget, and repeat
the prior-art and certificate checks. This theorem cannot be extrapolated
to those alternatives without additional proof.

\begin{thebibliography}{9}
\bibitem{bower} A. Bower, R. Evans, V. Luo, and S. J. Miller.
\emph{Coordinate sum and difference sets of $d$-dimensional modular
hyperbolas}. arXiv:1212.2930v1, submitted 12 December 2012.
Proposition 2.1, Lemma 2.3, Theorems 3.3 and 3.6.
\url{https://arxiv.org/abs/1212.2930v1}.

\bibitem{hittmeir} M. Hittmeir. \emph{Integer Factorization as Subset-Sum
Problem}. Journal of Number Theory 249 (2023), 93--118.
Reviewed version: arXiv:2205.10074v1, 20 May 2022; Sections 3, 4 and 7.
\url{https://arxiv.org/abs/2205.10074v1}.

\bibitem{hewitt} S. Hewitt. \emph{Empirical explorations of faster Fermat
factorisation, part 5}. Cambridge Clarion, 12 August 2022, updated
16 August 2022. Self-authored empirical study, section ``A systematic
approach to tuning the filter'' and appendix.
\url{https://www.cambridgeclarion.org/61.html}.
\end{thebibliography}
\end{document}
